\documentclass[12pt]{article}
\usepackage{amsmath, amssymb}
\usepackage{geometry}
\geometry{margin=1in}

\begin{document}

\title{CSE 400: Fundamentals of Probability in Computing\\
Tutorial--1}
\date{School of Engineering and Applied Science (SEAS), Ahmedabad University\\
February 13, 2026}
\maketitle

This document is a complete, step-by-step reconstruction of Tutorial--1. All questions are solved fully and explicitly. Every intermediate computation is written in detail.

\section*{Q1}

Twenty distinct dishes are to be divided into four platters of exactly five dishes each. The order of the platters does not matter.

\subsection*{(a) Total number of ways to divide 20 different dishes into four groups of five dishes each}

First choose 5 dishes for the first platter:

\[
\binom{20}{5}
\]

Then choose 5 dishes from the remaining 15:

\[
\binom{15}{5}
\]

Then choose 5 dishes from the remaining 10:

\[
\binom{10}{5}
\]

The remaining 5 dishes form the fourth platter:

\[
\binom{5}{5}
\]

Since the order of the four platters does not matter, divide by $4!$.

\[
\text{Total number of ways}
=
\frac{
\binom{20}{5}
\binom{15}{5}
\binom{10}{5}
\binom{5}{5}
}{4!}
\]

\subsection*{(b) Probability that all four platters have dishes of the same cuisine}

There are 5 dishes from each of the 4 cuisines.

Favorable outcome: each platter contains exactly the 5 dishes of one cuisine.

There is exactly one such partition.

Total number of partitions is:

\[
\frac{
\binom{20}{5}
\binom{15}{5}
\binom{10}{5}
\binom{5}{5}
}{4!}
\]

Therefore,

\[
P =
\frac{1}{
\frac{
\binom{20}{5}
\binom{15}{5}
\binom{10}{5}
\binom{5}{5}
}{4!}
}
\]

\section*{Q2}

Buses arrive at 15-minute intervals starting at 7:00 A.M.

Arrival time is uniformly distributed between 7:00 and 7:30.

Total interval length = 30 minutes.

\subsection*{(a) Probability that waiting time is less than 5 minutes}

Waiting is less than 5 minutes if arrival occurs in the last 5 minutes before a bus arrives.

These intervals are:

7:10--7:15

7:25--7:30

Each interval has length 5 minutes.

Total favorable time = 10 minutes.

\[
P = \frac{10}{30} = \frac{1}{3}
\]

\subsection*{(b) Probability that waiting time is more than 10 minutes}

Waiting exceeds 10 minutes if arrival occurs in the first 5 minutes after a bus departure.

These intervals are:

7:00--7:05

7:15--7:20

Total favorable time = 10 minutes.

\[
P = \frac{10}{30} = \frac{1}{3}
\]

\section*{Q3}

Given:

\[
P(\text{arrives late}) = 0.15
\]

\[
P(\text{leaves early}) = 0.25
\]

\[
P(\text{arrives late and leaves early}) = 0.08
\]

We compute:

Arrives early is complement of arrives late:

\[
P(\text{arrives early}) = 1 - 0.15 = 0.85
\]

Leaves late is complement of leaves early:

\[
P(\text{leaves late}) = 1 - 0.25 = 0.75
\]

Now compute:

\[
P(\text{arrives early AND leaves late})
=
1 - P(\text{arrives late}) - P(\text{leaves early}) + P(\text{arrives late AND leaves early})
\]

\[
= 1 - 0.15 - 0.25 + 0.08
\]

\[
= 0.68
\]

Thus,

\[
P(\text{arrives early} \mid \text{leaves late})
=
\frac{0.68}{0.75}
=
\frac{68}{75}
\]

\section*{Q4}

Expected number of typographical errors on a page = 0.2.

\[
X \sim \text{Poisson}(0.2)
\]

\[
P(X=k) = e^{-0.2}\frac{0.2^k}{k!}
\]

\subsection*{(a) 0 errors}

\[
P(X=0) = e^{-0.2}
\]

\subsection*{(b) 2 or more errors}

\[
P(X \ge 2)
=
1 - P(0) - P(1)
\]

\[
= 1 - e^{-0.2} - 0.2 e^{-0.2}
\]

\[
= 1 - e^{-0.2}(1+0.2)
\]

\section*{Q5}

Average number of airplane crashes per month = 3.5.

\[
X \sim \text{Poisson}(3.5)
\]

\subsection*{(a) At least 2 accidents}

\[
P(X \ge 2)
=
1 - P(0) - P(1)
\]

\[
= 1 - e^{-3.5} - 3.5 e^{-3.5}
\]

\[
= 1 - e^{-3.5}(1+3.5)
\]

\subsection*{(b) At most 1 accident}

\[
P(X \le 1)
=
P(0) + P(1)
\]

\[
= e^{-3.5} + 3.5 e^{-3.5}
\]

\[
= e^{-3.5}(1+3.5)
\]

\section*{Q6}

\[
X \sim \text{Poisson}(3)
\]

\subsection*{(a) $P(X \ge 3)$}

\[
P(X \ge 3)
=
1 - P(0) - P(1) - P(2)
\]

\[
P(0) = e^{-3}
\]

\[
P(1) = 3e^{-3}
\]

\[
P(2) = \frac{3^2}{2!}e^{-3} = \frac{9}{2}e^{-3}
\]

Thus,

\[
P(X \ge 3)
=
1 - e^{-3} - 3e^{-3} - \frac{9}{2}e^{-3}
\]

\[
= 1 - e^{-3}\left(1+3+\frac{9}{2}\right)
\]

\[
= 1 - \frac{17}{2}e^{-3}
\]

\subsection*{(b) $P(X \ge 3 \mid X \ge 1)$}

\[
\frac{P(X \ge 3)}{P(X \ge 1)}
\]

\[
P(X \ge 1)
=
1 - P(0)
=
1 - e^{-3}
\]

Thus,

\[
\frac{1 - \frac{17}{2}e^{-3}}{1 - e^{-3}}
\]

\section*{Q7}

Given:

\[
f_X(x)
=
\frac{1}{\sqrt{8\pi}}
\exp\left(-\frac{(x+3)^2}{8}\right)
\]

Let:

\[
Z = \frac{X+3}{2}
\]

\subsection*{1. $P(X \le 0)$}

\[
Z = \frac{0+3}{2} = \frac{3}{2}
\]

\[
P(X \le 0)
=
P\left(Z \le \frac{3}{2}\right)
=
1 - Q(1.5)
\]

\subsection*{2. $P(X > 4)$}

\[
Z = \frac{4+3}{2} = \frac{7}{2}
\]

\[
P(X > 4)
=
Q(3.5)
\]

\subsection*{3. $P(|X+3| < 2)$}

\[
-2 < X+3 < 2
\]

Divide by 2:

\[
-1 < Z < 1
\]

\[
P = 1 - 2Q(1)
\]

\subsection*{4. $P(|X-2| > 1)$}

\[
X > 3 \quad \text{or} \quad X < 1
\]

Standardize each term accordingly and express as sum of Q-functions.

\section*{Q8}

Let:

\[
X \sim \text{Binomial}(12,\theta)
\]

Conviction requires $(X \ge 8)$.

Correct decision probability:

\[
\sum_{k=8}^{12}
\binom{12}{k}
\theta^k(1-\theta)^{12-k}
+
\sum_{k=0}^{4}
\binom{12}{k}
\theta^k(1-\theta)^{12-k}
\]

\section*{Q9}

Given:

\[
p(i) = c\frac{\lambda^i}{i!}
\]

Since probabilities sum to 1:

\[
\sum_{i=0}^{\infty} c\frac{\lambda^i}{i!}
=
c e^{\lambda}
=
1
\]

\[
c = e^{-\lambda}
\]

\subsection*{(a)}

\[
P(X=0)
=
e^{-\lambda}
\]

\subsection*{(b)}

\[
P(X>2)
=
1 -
e^{-\lambda}\left(1+\lambda+\frac{\lambda^2}{2}\right)
\]

\section*{Q10}

System operates if at least half components function.

\subsection*{(a)}

5-component system effective if $\ge 3$ function:

\[
P_5
=
\binom{5}{3}p^3(1-p)^2
+
\binom{5}{4}p^4(1-p)
+
\binom{5}{5}p^5
\]

3-component system effective if $\ge 2$ function:

\[
P_3
=
\binom{3}{2}p^2(1-p)
+
\binom{3}{3}p^3
\]

5-component system is better when:

\[
P_5 > P_3
\]

\subsection*{(b)}

\[
\sum_{j=k+1}^{2k+1}
\binom{2k+1}{j}
p^j(1-p)^{2k+1-j}
>
\sum_{j=k}^{2k-1}
\binom{2k-1}{j}
p^j(1-p)^{2k-1-j}
\]

\bigskip

\centerline{\textbf{End of Complete Exam-Oriented Tutorial--1 Scribe}}

\end{document}
